% 张三慧《学习辅导与习题解答》PDF S 独立选答，卷1第5章；未取得V1主教材独立交叉核验。
% 仅录“思考题选答”中实际明确列出的原编号；其余不推断，exam_ready=false。
\section*{第5章\quad 刚体的转动：辅导书已列思考题独立解析}
\paragraph{思考题 5.3}
长杆可使演员—杆系统绕可能的翻倒转轴具有更大转动惯量 \(I=\int r_\perp^2dm\)。同一扰动力矩 \(\tau\) 下，角加速度 \(\alpha=\tau/I\) 变小，给予演员更多调整时间。杆的质量分布还能改变质心和恢复力矩，但单凭大转动惯量不能在无主动控制的所有姿态下保证稳定。

\paragraph{思考题 5.6}
冰面摩擦力矩若可忽略，关于竖直轴角动量 \(L_z=I_z\omega\) 近似守恒。收臂减小质量到轴的平方距离，令转动惯量由 \(I_1\) 降至 \(I_2\)，故 \(\boxed{\omega_2/\omega_1=I_1/I_2>1}\)。转动动能 \(K=L_z^2/(2I)\) 因而增加，其增加来自运动员克服离心效应做的肌肉功，不违反能量守恒。

\paragraph{思考题 5.10}
旋转盘具有角动量 \(\boldsymbol L\approx I_3\omega_s\hat{\boldsymbol e}_3\)。重力对支点产生力矩 \(\boldsymbol\tau=\boldsymbol r_{\rm cm}\times m\boldsymbol g\)，满足 \(\dot{\boldsymbol L}=\boldsymbol\tau\)，主要改变角动量\emph{方向}，产生陀螺进动，而非必然使盘立刻倾倒。简单慢进动近似 \(\Omega_p\simeq mgl/(I_3\omega_s)\)；当初始条件未恰好满足规则进动或转速不足，还会出现周期性章动。旋转本身不产生永久稳定的“神秘支撑力”。
